\chapter{The Sense of Consonance}

\epigraph{\emph{Spiritual homeostasis.}}{}

Chapter 8 proposed that some forms of natural harmonic organization can function as attractors of informational efficiency. That proposal remains incomplete unless it can also explain why such organization should matter to living systems as something more than an abstract formal convenience. If efficient recurrent regimes lower corrective burden while preserving usable difference, then the next question is biological before it is aesthetic: can such regimes become detectable, preferable, or regulatory for organisms that must remain viable under conditions of noise, variation, and constraint? Chapter 9 takes up that question by asking how consonance can begin to operate as a sense.

The chapter therefore moves one step away from abstract dynamics and one step closer to embodied orientation. Its argument is not that consonance has already been proven to heal, to calm, or to reveal some hidden truth of the subject. Nor does it claim that any pleasant sonic experience should be re-described as harmony in the strong HIT sense. The narrower proposal is that under some conditions consonant organization may function as a candidate orientation regime: a way in which living systems register lower-friction coordination, reduced overload, or greater coherent relation among ongoing processes. That proposal will remain conditional throughout. The chapter is concerned with plausibility, articulation, and convergence, not with a completed therapeutic or mechanistic proof.

\section{From efficiency to living orientation}

The transition from efficiency to orientation becomes possible once efficiency is understood in biological rather than merely engineering terms. Organisms do not persist by minimizing expenditure in the abstract. They persist by maintaining viable organization while continuously negotiating perturbation. In that setting, a regime that lowers corrective burden may matter because it helps preserve coherence without requiring constant large-scale readjustment. Consonance becomes theoretically relevant here not as a musical category alone, but as one candidate family of relations through which such low-conflict coordination can be sensed, sustained, or re-entered.

This is why the chapter prefers the expression homeostatic orientation to the stronger claim of homeostasis already achieved. A living system rarely inhabits perfect equilibrium. It continually approaches, loses, and partially regains coordinated states across multiple timescales. Oscillatory physiology provides a natural archive for this point. Glass and Mackey show that biological rhythms remain viable through flexible coordination among partially independent cycles rather than through rigid uniformity \parencite{glass_1988}. Acebron and colleagues make an analogous point in coupled-oscillator language: partial synchronization, chimera-like distributions, and selective locking are not failures of order but common regimes of organized persistence \parencite{acebron_2005}. Consonance matters in that context because it can be read as one class of relations that tends to support return, re-entry, and lower-conflict coupling without abolishing difference.

Recent neurophysiological work gives this intuition a more literal anchor. Kawai found slow harmonic couplings among cardiorespiratory and brainstem rhythms in vivo, and Zheng and colleagues showed that near-harmonic alpha-theta relations modulate attentional steadiness in humans \parencite{kawai_2023, zheng_2025}. In both cases, orientation is not a metaphor added afterward. The operative state of the system changes with the quality of the relation it can sustain.

Porges's polyvagal work sharpens the regulatory dimension of that claim. If autonomic organization is sensitive to patterns of safety, modulation, and environmental predictability, then structured sensory input may matter not only because it is detected, but because it participates in the conditions under which a system shifts among defensive, mobilized, and socially regulated states \parencite{porges_2011, porges_2017}. Levin's work on bioelectric organization extends the same intuition beyond the nervous system narrowly conceived: living form depends on coordinated ratios, gradients, and patterned signaling rather than on isolated local events alone \parencite{levin_2021}. Taken together, these literatures do not establish that harmonic structure directly causes well-being. They support the more disciplined inference that some relational environments may function as orienting conditions for biological regulation, and that consonant patterning is a plausible candidate among them.

Under this reading, the value of consonance is not exhausted by pleasantness. A relation can matter because it furnishes a recoverable reference. The relevant image is not mystical harmony but reorientable coherence: the possibility that a system drifting across perturbation can regain organization more readily when the field it encounters is structured by legible, low-conflict relations. Later chapters will ask whether the Harmonic Beacon, the experiential and acoustic-physiological branch developed in Chapter 12, can operationalize something like that reference in practice. At the present stage the claim remains more general. If natural harmonic regimes can lower corrective burden, then organisms may be able to register them as conditions of easier coordination rather than as neutral stimuli among others.

\section{Consonance, affect, and predictive ease}

Once consonance is treated as a candidate orienting regime, its affective dimension becomes easier to describe without romanticizing it. What a system senses as ease, relief, or attraction need not be interpreted as mere subjective decoration added after the fact to a formally efficient arrangement. Affective valence may be one way in which a system registers the difference between higher-friction and lower-friction states of organization. In humans, the relevant archive includes reward, expectancy, tension, release, and autonomic modulation rather than musical preference in the narrow cultural sense alone.

Blood and Zatorre's work on intensely pleasurable musical experience provides an early and still powerful anchor for this point. Their results show that peak responses to music involve circuitry also implicated in reward and motivation, including the nucleus accumbens and related dopaminergic pathways \parencite{blood_2001}. Ferreri and colleagues strengthen the bridge by demonstrating that pharmacological modulation of dopamine alters musical pleasure in measurable ways \parencite{ferreri_2019}. Chanda and Levitin, Koelsch, and Panksepp and Bernatzky each contribute to the same picture from slightly different angles: musical patterning is capable of modulating affect, reward, and emotion-laden processing in ways that are neither trivial nor reducible to cultural convention alone \parencite{chanda_2013, koelsch_2014, panksepp_2002}. None of this yet isolates natural harmonicity as the decisive variable, but it does show that patterned sonic organization can penetrate deeply into systems of valuation and regulation.

Predictive frameworks clarify why that penetration may matter. If the brain persistently anticipates the structure of incoming signals, then patterns that stabilize expectation without eliminating discriminability may be experienced as easier to inhabit than patterns that impose continual corrective work. Friston's framework is useful here not because it turns consonance into a solved equation, but because it provides a disciplined vocabulary for saying that lower surprise and lower corrective burden can be affectively legible \parencite{friston_2010}. The lesson carried forward from Chapter 8 remains decisive: the target regime is not total order. A perfectly repetitive tone with no usable variation would be predictively simple and experientially thin. The plausible sweet spot is narrower and more interesting: enough recurrence to reduce friction, enough variation to preserve attention, difference, and response.

Large and Almonte's resonance-based account of tonal perception helps bridge the formal and affective sides of the problem \parencite{large_2012}. Tonal relations are not simply decoded as abstract ratios; they recruit dynamic resonance in the perceiving system. In that sense, consonance may be felt because it is not only represented but enacted as easier coordination within the listener. Restorative and stress-related literatures gesture toward similar possibilities at the edge of the argument. Thoma and colleagues, de Witte and colleagues, and Bradt and colleagues all indicate that patterned musical exposure can in some contexts reduce stress markers or support better regulation, even though those findings concern heterogeneous interventions and do not isolate natural harmony in the strict sense \parencite{thoma, de_witte, bradt_2021}. The prudent conclusion is therefore not that consonance has been shown to heal, but that affect may be one way living systems register relative predictive ease and reduced corrective demand.

\section{Disharmony, overload, and estrangement}

If consonance can be approached as orientation, then dissonance or inharmonicity cannot be treated merely as its aesthetic opposite. They must also be considered as possible names for states of overload, unresolved competition, masking, or estrangement between a system's current organization and the field in which it moves. That claim, however, requires restraint. HIT should not collapse into a moralized contrast between pure nature and corrupt city, nor into the fantasy that all technological environments are intrinsically pathological because they are dense, amplified, or noisy. The relevant distinction is not rural virtue versus urban vice. It is legible coordination versus sustained overload under given conditions.

Acoustic ecology offers one disciplined way to formulate the problem. Krause's work on the acoustic niche and Buxton and colleagues' work on the restorative value of natural soundscapes suggest that environments differ not only in loudness, but in how signals are distributed, masked, and made available for interpretation \parencite{krause_2013, buxton_2021}. A field saturated by poorly differentiated or relentlessly competing signals can force a receiving system into continual filtering, defensive gating, or missed discrimination. In that sense, estrangement may be understood less as alienation in the purely philosophical register and more as a mismatch between the organism's capacities for coordination and the structure of the environment it must inhabit. The cost is then not merely annoyance. It is sustained corrective labor.

Koelsch's work on emotionally salient auditory processing helps keep this point tied to physiology rather than cultural lament \parencite{koelsch_2014}. Environments that overtax discrimination, regulation, or anticipation can recruit stress, vigilance, and defensive organization. HIT does not need to claim that all inharmonic or noisy environments are harmful, because that would both overstate the case and erase the productive uses of tension, rupture, and novelty. The narrower proposal is that chronically overloaded environments may distance a system from lower-conflict regimes of coordination, while consonant patterning may sometimes provide one route of reorientation. In that precise sense, disharmony names not a moral impurity but a relative increase in friction between ongoing organization and the conditions required to sustain it.

This distinction also prevents a conceptual mistake. If consonance is always identified with comfort and dissonance with damage, the theory becomes sentimental. Many living systems need perturbation, challenge, and contrast in order to remain adaptive. What matters is whether tension can be metabolized into a new coherence or whether it accumulates as unresolved burden. The chapter's use of estrangement should therefore be read as structural rather than moral: a condition in which legible relation becomes harder to sustain, not a universal diagnosis of the modern world.

\section{Harmonic pareidolia and human polyphony}

The argument becomes more recognizably human when it turns from regulation alone to the search for pattern. Perception is not a passive reception of finished objects. It is an active organization of signals into structures that can be recognized, anticipated, and acted upon. In that setting, pareidolia is not a curiosity at the margins of cognition. It is a reminder that living systems prefer interpretable organization to undifferentiated flux. The question for HIT is whether some forms of harmonic patterning participate in that search more deeply than others.

Nees and Phillips show that listeners often impose meaningful organization on ambiguous auditory material, revealing that auditory perception is structured by expectation, grouping, and pattern completion rather than by bare stimulus registration alone \parencite{nees_phillips}. Nikolsky's historical and comparative work on scales and sonic organization suggests that human tonal systems repeatedly stabilize around relations that are cognitively tractable, socially iterable, and perceptually salient rather than arbitrarily distributed across the possible frequency continuum \parencite{nikolsky}. Schwartz and colleagues add an ecological layer by showing that speech itself contains non-random statistical regularities that connect vocal production, perception, and harmonic structure in ways that are unlikely to be accidental \parencite{schwartz_2003}. These literatures do not imply that the brain simply \enquote{finds} natural harmony waiting fully formed in the world. They suggest instead that systems of perception are biased toward extracting recurrent, usable structure from complex signals, and that harmonic organization may be one privileged case of such extraction.

This is one place where the intuition that humans are polyrhythmic and polyphonic can be translated into academic prose without being flattened into banality. Human vocalization is never a single abstract note. The voice emits harmonic complexes whose prosodic contour, overtone structure, timbral grain, breath pattern, and bodily tension all participate in how affect and intention are heard. Koelsch and related work on emotional prosody indicate that the sonic envelope of a voice carries structured information about state, disposition, and interpersonal orientation \parencite{koelsch_2014}. Panksepp and Bernatzky similarly help frame music and vocal patterning as routes into deep affective organization rather than superficial ornament \parencite{panksepp_2002}. To say that humans are polyphonic, then, is not to indulge metaphor. It is to say that emotional life is continuously emitted, perceived, and coordinated through multilayered sonic constellations rather than through isolated tones alone.

This matters for HIT because consonance may operate not only as a property of external sound, but as one schema through which internal and interpersonal organization is read. A stable interval, a coherent timbral field, or a voice whose overtones and rhythm remain legible can all become cues for relation. Conversely, ambiguous or unstable patterning can intensify interpretive labor. Harmonic pareidolia is therefore not just a tendency to hear order where none exists. It is also a sign that organisms live by constructing actionable pattern from noisy fields. That observation prepares the chapter's next step, where consonance is reformulated more precisely as a biologically enacted relation of structural coupling. Only from that bridge can the final step toward psychic orientation be taken without equivocation.

\section{Consonance as structural coupling}

At this point the chapter can name the biological mechanism it has so far approached only indirectly. From an autopoietic and enactive perspective, the sense of consonance is not best described as the passive registration of an external ratio that later acquires value. It is better described as a relation of structural coupling in which a living system encounters perturbations that are more or less compatible with its own organization. Maturana's rejection of information transfer is decisive here. Organisms do not receive information from the environment as if messages entered a neutral channel. They are perturbed, and the significance of the perturbation depends on the structure and history of the receiving system \parencite{maturana_1970, maturana_varela_1980}. Under that description, consonance becomes legible as a regime in which the perturbing field can be integrated with lower conflict than nearby alternatives. The relevant point is not that harmony is injected into the organism from outside, but that some relations prove more compatible with ongoing organization than others.

This is one reason the enactive vocabulary is so useful at this juncture. In \emph{The Embodied Mind}, cognition is not defined as internal representation of a pregiven world, but as effective action within a domain of coupling \parencite{varela_thompson_rosch_1991}. The sense of consonance can therefore be redescribed as enacted compatibility rather than as contemplative recognition alone. A system does not first decode an external relation and only afterward respond to it. It finds itself more or less able to continue coherently within a field of perturbation. That is also why the chapter's earlier appeal to homeostatic orientation can now be stated more precisely. The issue is not whether consonance is always pleasant. The issue is whether some patterned relations allow a precarious system to preserve organization with less unresolved correction than others do.

Di Paolo's notion of precarious autonomy gives this account its normative edge. The consonant matters because a living system can fail, fatigue, overload, or disintegrate in its coordination with the environment \parencite{dipaolo_2005}. In that sense, inharmonicity threatens organization not simply because it sounds unpleasant, but because it may require more filtering, more compensation, or more unresolved adaptive labor. The claim should remain narrow. HIT is not saying that organisms need harmonic environments in order to survive in every case, nor that urban or noisy contexts are biologically illegitimate. It is saying that the language of structural coupling provides a more precise biological vocabulary for the difference between supportive and burdensome perturbation than the looser language of preference alone. Even the Harmonic Beacon, introduced fully in Chapter 12 as the program's experiential and acoustic-physiological branch, should be described in those terms: not as a device that transmits harmony into the body, but as a controlled field of perturbation to which differently structured bodies may couple differently.

In \emph{La tecnica digitalizada}, Fern\'andez M\'endez describes the pilot--all-terrain-vehicle relation as a striking example of what such coupling can look like outside abstract system language. There, successful operation depends on the subjective incorporation of the vehicle into the body's own functional schema, such that machine and operator no longer confront one another as merely external terms \parencite{fernandez_2021b}. Something structurally similar is at stake in the present account of consonance. The relevant issue is not identity or fusion, but a form of enacted compatibility in which the system can extend its organized action through a field that would otherwise remain external, noisy, or resistant.

Varela's work on large-scale neural integration suggests one further, still hypothetical bridge. If distributed neural coherence is indeed achieved in part through transient phase synchronization among resonant ensembles, as the Brainweb literature argues, then a perceptual field with stronger harmonic coherence may facilitate rather than hinder the formation of integrative neural relations \parencite{varela_2001_brainweb, rodriguez_1999}. That stronger reading remains open and must stay open. The chapter is not entitled to claim that Beacon already produces such integration or that consonance has been shown to optimize neural coupling in vivo. What it can say is more modest and more important: these traditions help explain why consonance might matter as an enacted sense of compatibility rather than only as an acoustic or cultural category.

At that point the stronger expression \emph{spiritual homeostasis} becomes available. By spiritual homeostasis, HIT refers to the recurrent capacity of a living system to recover and sustain patterns of resonance, orientation, and coherence across biological, affective, symbolic, and ecological levels. It does not name final equilibrium. It names a livable order regained through reorganization, one in which dispersion, overload, or dissonant pressure are brought back into a form the system can inhabit without losing differentiation. In that sense, consonance matters not only as an acoustic relation or a regulatory aid, but as a possible operator of reorientation across coupled layers of experience.

That account, however, does not exhaust the problem of orientation. Structural coupling explains how a living system preserves or loses congruence with its environment, but it does not yet describe the specifically psychic forms through which tension, symbolization, conflict, and reorganization are lived in human experience. If HIT is to move from enacted compatibility to psychic orientation without collapsing one register into the other, it needs a second vocabulary. This is the point at which Jung becomes useful.

\section{Consonance as orienting reference}

The chapter's final step requires a change of register, but not a break with the earlier argument. If consonance can function as an orienting condition at biological and affective levels, then psychology offers a plausible next question: can it also serve as a figure for how the psyche is guided toward greater coherence without ever being fully transparent to itself? The analytic tradition of Jung is useful here because it names an organizing factor that is neither reducible to the conscious ego nor exhausted by the unconscious taken as a mere reservoir. The Self is presented instead as a principle of orientation that becomes legible through the process of individuation and through the tensions that gradually organize the relation between conscious and unconscious life \parencite{jung_1959, vonfranz_1964}.

This matters to HIT because the Self, read cautiously, resembles the kind of orienting reference the theory has been describing elsewhere without sharing its ontology or substrate. It is not a fixed entity hidden behind appearances. It is a relational center that manifests through processes of integration, conflict, symbolic emergence, and partial reorganization. Jung's account of the transcendent function sharpens the analogy. The function does not abolish opposition by choosing one pole over another. It produces a third position in which tension is worked through into a form that neither side could generate alone \parencite{jung_1916}. HIT can therefore borrow a disciplined analogy: natural harmonic reference may orient a system toward coherence not by erasing complexity, but by giving conflicting tendencies a relation in which they can be partially integrated.

Campbell's monomyth gives the same logic a narrative morphology \parencite{campbell_1949}. Departure, descent, confrontation, and return are not relevant here because the manuscript wishes to mythologize consonance, but because they name a repeatable structure in which disorientation and reorientation are inseparable. A system loses its previous equilibrium, enters a zone of intensified tension, and re-emerges with a reorganized relation to itself and its world. In that narrow sense, the hero's path resembles the movement already described in biological and affective terms: not permanent equilibrium, but recurrent passage through instability toward renewed coherence. Later, Chapter 12 will ask whether the Harmonic Beacon, introduced there as the experiential branch of the program, can help stage or support conditions under which such reorientation becomes experientially investigable. At this stage the claim is smaller. The consonant may name not only lower-friction patterning, but also an orienting function for psychic life.

The epistemic caution here must remain explicit. HIT is not claiming that the Self is literally a harmonic ratio, that individuation can be reduced to acoustics, or that psychic life is transparently measurable by consonance alone. The analogy is programmatic and methodological, not ontological. It extends the argument of Chapter 7 rather than contradicting it. The psyche, like any object that matters to this project, exceeds complete formal capture and appears through mediated effects, symbolic formations, and bodily events rather than through direct possession of the Real. That is precisely why the Jungian reference is useful. It provides a way to think orientation without demanding full transparency. Consonance can therefore be proposed as a possible psychic sense not because it reveals essence, but because it may help describe how systems grope toward coherence across registers that remain partially opaque to formalization.

Chapter 9 has moved the argument as far as it can go without confronting a missing problem in the theory itself. If consonance can be understood as a candidate orientation function across regulation, affective ease, environmental fit, and psychic reorganization, then one question remains before the program can turn to its probes: how does a harmonically organized system become readable without having its organization destroyed by the very act of reading? The next chapter takes up that question as a problem of activation and retrieval. Only after that does the book turn to Phideus and the Harmonic Beacon as distinct experimental branches of the same program.
